\documentclass[12pt]{article}
\usepackage{bbm}
\usepackage{fancyhdr}
\usepackage{titling}
\usepackage{enumitem}
\usepackage{amsmath}
\usepackage{dsfont}
\usepackage{amsfonts}
\usepackage{amssymb}
\usepackage{amsthm}
\usepackage{geometry}
\usepackage{tikz}
\usepackage{algorithm}
\usepackage{algpseudocode}
\geometry{verbose,tmargin=1in,bmargin=1in,lmargin=1in,rmargin=1in}
\renewcommand{\vec}[1]{\mathbf{#1}}
\newcommand{\dec}{\text{dec}}
\newcommand{\divi}{\mathbin{\textnormal{div}}}
\title{Homework 1}
\author{Satvik Prasad}
\date{}
\newtheorem{lemma}{Lemma}[section]
\fancypagestyle{plain}{%
  \fancyhf{}
  \fancyhead[L]{\theauthor}
  \fancyhead[R]{\thedate}
}
\pagestyle{plain}
\makeatletter
\def\@maketitle{%
  \newpage
  \null
  \vskip 1em%
  \begin{center}%
    \let \footnote \thanks
    {\LARGE \@title \par}%
    \vskip 1em%
  \end{center}%
  \par
\vskip 1em}
\makeatother
\usepackage{cmbright}

\newcommand{\mat}[1]{\mathbf{#1}}

% ---------- Probability theory ----------
% Core operators (optional argument-free; use as \E{X}, \Var{X}, \Cov{X}{Y})
\DeclareMathOperator{\EOp}{\mathbb{E}}
\DeclareMathOperator{\VarOp}{Var}
\DeclareMathOperator{\CovOp}{Cov}
\DeclareMathOperator{\CorrOp}{Corr}

\renewcommand{\P}[1]{\mathbb{P}\left(#1\right)}        % \P{A}
\newcommand{\E}[1]{\EOp\left[#1\right]}                % \E{X}
\newcommand{\Var}[1]{\VarOp\left(#1\right)}            % \Var{X}
\newcommand{\Cov}[2]{\CovOp\left(#1,\,#2\right)}       % \Cov{X}{Y}
\newcommand{\Corr}[2]{\CorrOp\left(#1,\,#2\right)}     % \Corr{X}{Y}
\newcommand{\SD}[1]{\operatorname{SD}\left(#1\right)}  % \SD{X}

% Subscripted versions, e.g. \Esub{\theta}{X}, \Psub{\theta}{A}
\newcommand{\Esub}[2]{\EOp_{#1}\left[#2\right]}
\newcommand{\Psub}[2]{\mathbb{P}_{#1}\left(#2\right)}

% Conditioning, indicators, independence
\newcommand{\given}{\,\middle|\,}                      % \E{X \given Y}
\newcommand{\ind}[1]{\mathbbm{1}\left\{#1\right\}}     % \ind{X > 0}
\newcommand{\indep}{\perp\!\!\!\perp}                  % X \indep Y
\newcommand{\iid}{\overset{\text{iid}}{\sim}}          % X_i \iid \Normal(0,1)

% Convergence
\newcommand{\convp}{\xrightarrow{\;p\;}}               % in probability
\newcommand{\convd}{\xrightarrow{\;d\;}}               % in distribution
\newcommand{\convas}{\xrightarrow{\text{a.s.}}}        % almost surely
\newcommand{\convLp}[1]{\xrightarrow{L^{#1}}}          % in L^p

% Probability spaces
\newcommand{\F}{\mathcal{F}}                           % sigma-algebra
\newcommand{\B}{\mathcal{B}}                           % Borel sigma-algebra
\newcommand{\pspace}{(\Omega, \F, \mathbb{P})}

% Common distributions
\newcommand{\Normal}{\mathcal{N}}
\newcommand{\Bern}{\operatorname{Bernoulli}}
\newcommand{\Bin}{\operatorname{Binomial}}
\newcommand{\Pois}{\operatorname{Poisson}}
\newcommand{\Unif}{\operatorname{Uniform}}
\newcommand{\Exp}{\operatorname{Exp}}
\newcommand{\Geom}{\operatorname{Geometric}}
\newcommand{\Gam}{\operatorname{Gamma}}
\newcommand{\Beta}{\operatorname{Beta}}

% Number sets
\newcommand{\R}{\mathbb{R}}
\newcommand{\N}{\mathbb{N}}
\newcommand{\Z}{\mathbb{Z}}


\begin{document}
\begin{enumerate}
  \item
    \begin{enumerate}
      \item Let us prove the forward direction first. Consider the
        case where we have a subspace $U \subseteq V$ with $U^{0} =
        \{0\}.$ Hence, the only functional $\varphi \in V'$ is the
        zero map, $0.$ For the sake of contradiction, assume $U \neq
        V.$ Then, we have some element $u_0 \in V \setminus \{0\}.$
        Hence, we can construct a linear map $\varphi$ such that
        $$
        u \neq \mu u_0,\,  \forall \mu \in \mathbb F \Rightarrow
        \varphi(u) = 0,
        $$
        and $\varphi(u_0) = \lambda$ for some scalar $\lambda \neq
        0.$ This map $\varphi \in U^{0},$ but is not the zero map,
        meaning $U^{0} \neq \{0\}.$ Hence, by contradiction $U = V.$

        Conversely, assume we have $U = V.$ Then, for all $\varphi
        \in U^{0},$  $\varphi(v) = 0$ for all $v \in V,$ since $V =
        U.$ Therefore, the only element in $U^{0}$ is the zero map,
        $0$, meaning $U^{0} = \{0\}.$

      \item Assume that we have $U^{0} = V'.$ Then, assume that $U
        \neq \{0\},$ i.e., there is some vector $0\neq u_0 \in U.$
        Then, any map $\varphi \in V'$ that satisfies $\varphi(u_0) =
        0$ is not in $U^{0}.$ Since such a map always exists, $U^{0}
        \neq V',$ since we have some $\varphi \in V'$ and $\varphi
        \not \in U^{0}.$ Therefore, by contradiction, $U = \{0\}.$

        Conversely, assume $U = \{0\}.$ Then, any map $\varphi \in
        V'$ satisfies $\varphi(0)= 0$ since $\varphi$ is linear,
        meaning $V \subseteq U^{0}.$ Since every element in $U^{0}$
        is linear, we have $U^{0} \subseteq V,$ and hence by dual set
        containment $U^{0} = V'.$
    \end{enumerate}

  \item
    \begin{enumerate}
      \item Assume that $T: V \to W$ is injective, meaning $\dim
        \text{null}\, T = 0,$ or equivalently, $\dim V = \dim
        \text{range}\, T.$ Then, $T'(\varphi) = \varphi \circ T$ for
        all $\varphi \in W'.$ Then, consider any map $\psi \in V'$.
        Let $V$ have basis $v_1, v_2, \ldots, v_n.$ Since $T$ is
        injective, $T(v_1), T(v_2), \ldots, T(v_n)$ are also linearly
        independent, and can be extended to a basis of $W.$ Construct
        such a basis,
        $$
        T(v_1), T(v_2), \ldots,T(v_n), w_1, \ldots, w_m
        $$
        for $m \geq 0.$ Then, the map $\varphi$ satisfying
        $$
        \varphi(T(v_k)) = \psi (v_k)
        $$
        and $\varphi(w_k) = 0$ for all such valid indices $k$
        satisfies $T'(\varphi) = \psi,$ meaning $T'$ is surjective.

      \item Assume that $T: V \to W$ is surjective. Then, assume we
        have $\varphi, \psi \in W'$ such that $T(\varphi) = T(\psi),$
        or equivalently,
        $$
        \varphi \circ T = \psi \circ T.
        $$
        Hence, $$
        \varphi(T(v)) = \psi(T(v))
        $$
        for all $v \in V.$ Since $T$ is surjective, $\text{range}\, T
        = W,$ meaning for all $w \in W$ there is some $v_0$ such that
        $T(v_0) = w.$ Hence, for all $w,$ $\varphi(w) = \psi(w),$
        meaning $\varphi = \psi$ as required, since functionals are
        unique. Thus, $T'$ is injective, since $$T' \varphi = T '
        \psi \Rightarrow \varphi = \psi.$$
    \end{enumerate}

  \item
    \begin{enumerate}
      \item We do this via double set containment. First, consider
        some element $\varphi \in (U + W)^{0}.$ Thus, for any $x \in
        V$ such that there exists $u \in U, w \in W$ with $x = u +
        w,$ we have $$
        \varphi(x) = 0 = \varphi(u) + \varphi(w).
        $$
        Since for all $u \in U,$ $w \in W,$ $x = u + w \in (U +
        W)^{0},$ $\varphi(u) + \varphi(w) = 0$ for all such $u, w.$
        Fixing $w = 0,$ we get $\varphi(u) = 0,$ and fixing $u = 0,$
        we get $\varphi(w) = 0,$ meaning $\varphi$ must send all
        elements of $U$ to $0,$ and all elements of $W$ to $0.$
        Hence, $\varphi$ is in both $U^{0}$ and $W^{0},$ meaning
        $$
        x \in U^{0} \cap W^{0},
        $$
        i.e., $(U + W)^{0} \subseteq U^{0} \cap W^{0}.$

        Conversely, consider $\varphi \in U^{0} \cap W^{0}.$ Then,
        for all $u \in U$ and $w \in W,$ $\varphi(u + w) = \varphi(u)
        + \varphi(w) = 0.$ Hence, for all $x \in U + W, $ $\varphi
        (x) = 0,$ meaning $\varphi \in (U + W)^{0}.$ Thus we know
        $U^{0} \cap W^{0} \subseteq (U + W)^{0},$ and thus $(U +
        W)^{0} = U^{0} \cap W^{0}.$

      \item If we have $\varphi \in U^{0} +W^{0},$ then $\varphi = t
        + s$ for $t \in U^{0}$ and $s \in W^{0}.$ Hence, for any $v
        \in U \cap W,$ $\varphi(v) = t(v) + s(v) = 0,$ i.e., $\varphi
        \in (U \cap W)^{0}.$ Therefore, we have $U^{0} + W^{0}
        \subseteq (U \cap W)^{0}.$

        Notice that
        \begin{align*}
          \dim(U^{0} + W^{0}) &= \dim U^{0} + \dim W^{0} - \dim U^{0}
          \cap W^{0} \\
          &= 2\dim V - \dim U - \dim W-  \dim (U + W)^{0} \\
          &= 2 \dim V - \dim U - \dim W - V + \dim (U + W) \\
          &= \dim V - \dim U - \dim W + \dim U + \dim W - \dim U \cap W \\
          &= \dim V - \dim U \cap W = \dim (U \cap W)^{0}.
        \end{align*}

        Since these vector spaces have the same dimension, and one is
        a subset of the other, we have
        $$
        U^{0} + W^{0} = (U \cap W)^{0}.
        $$
    \end{enumerate}

  \item Let $V'' = (V')'.$ Consider $
    \Lambda : V \to V''
    $ with $(\Lambda v)(\varphi) = \varphi(v).$
    \begin{enumerate}
      \item Consider $\Lambda(u + w).$ By definition, we have
        $\Lambda(u + w)(\varphi) = \varphi(u + w) = \varphi(u) +
        \varphi(w)$ by the linearity of $\varphi.$ Since $\varphi(u)
        = (\Lambda u)(\varphi),$ and $\varphi(w) = (\Lambda
        w)(\varphi),$ we have
        $$
        (\Lambda(u + w))(\varphi) = (\Lambda u)(\varphi) + (\Lambda
        w)(\varphi) = (\Lambda u + \Lambda w)(\varphi).
        $$
        Since this holds for all $\varphi \in V',$ we get $\Lambda (v
        + w) = \Lambda v + \Lambda w.$

        Similarly, consider $\Lambda(\lambda v).$ For all $\varphi, $ we have
        $$
        (\Lambda(\lambda v))(\varphi) =\varphi(\lambda v) = \lambda
        \varphi(v) = \lambda (\Lambda v)(\varphi).
        $$ Since $\Lambda v \in V'',$ which is linear,  we get
        $$
        (\Lambda(\lambda v))(\varphi) = (\lambda \Lambda v)(\varphi)
        $$
        for all $\varphi,$ meaning
        $$
        \Lambda(\lambda v) = \lambda \Lambda(v).
        $$
        Thus, $\Lambda$ is lienar.

      \item Assume we have some $T \in \mathcal L(V).$ Then, consider
        $T'' \circ \Lambda.$ By definition, $T': V' \to V'$ satisfies
        $$
        T'(\varphi) = \varphi \circ T,
        $$
        and hence $T''(\psi) = \psi \circ T'.$ Then,  for any $v \in V,$
        $$
        (T'' \circ \Lambda)(v) = T''(\Lambda(v)) = \Lambda (v) \circ  T'.
        $$
        Therefore,  by the construction of $\Lambda,$ we have
        $$
        (\Lambda(v) \circ T')(\varphi) = \Lambda(v)(T'(\varphi)) =
        \Lambda(v)(\varphi \circ T) = (\varphi \circ T)(v) =
        \varphi(T(v)) = \Lambda(T(v)).
        $$
        Hence, for all $v \in V,$ we get
        $$
        (T'' \circ \Lambda)(v) = (\Lambda \circ T)(v),
        $$
        and thus $T'' \circ \Lambda = \Lambda \circ T.$

      \item Let us look at the null space of $\Lambda.$ Assume for an
        arbitrary $v,$
        $$
        \Lambda(v) = 0 \in V''.
        $$
        Then, for all $\varphi \in V',$ we have
        $$
        (\Lambda(v))(\varphi) = 0 = \varphi(v),
        $$
        meaning $v \in \ker \varphi$ for all such $\varphi.$ Then,
        for the sake of contradiction assume $v \neq 0.$ Then, all
        linear functionals $\varphi \in V'$ must have $\varphi(v) =0.$

        Now, consider the case where $\dim V > 0$. Since $V$ is
        finite dimensional, and $\dim V \neq 0,$ extend $v$ into a
        basis $v, v_1, \ldots, v_n$ with
        $$
        \varphi(v) = 1, \varphi(v_k) = 0.
        $$
        This is a contradiction, since $\varphi(v) = 0$ for all $v
        \in V,$ and thus $v = 0.$ Thus, $\ker \Lambda = 0,$ and by rank-nullity,
        $$
        \dim V'' = \dim V= 0 + \dim \text{range} \ \Lambda = \dim
        \text{range} \ \Lambda.
        $$
        Thus, $\Lambda$ is an isomorphism.

        Otherwise, if $\dim V = 0,$ then $\dim V'' = 0  = \dim \ker
        \Lambda + \dim \text{range}\  \Lambda.$ Since $\dim \cdot
        \geq 0,$ we have $\dim \text{range} \ \Lambda = 0 = \dim V,$
        meaning $\Lambda$ is till an isomorphism.
    \end{enumerate}

  \item Let $\pi'$ be the dual map of $\pi.$
    \begin{enumerate}
      \item Assume that we have $\pi'(\varphi) = \pi'(\psi)$ for
        $\varphi, \psi \in (V/U)'.$ Hence, we have
        $$
        (\varphi \circ \pi)(v) = (\psi \circ \pi)(v)
        $$
        for all $v \in V.$ Equivalently, we get
        $$
        \varphi(v + U) = \psi(v + U)
        $$
        for all $v \in V.$ Since $\varphi = \psi$ over all of their
        domain $V/U,$ we get $\varphi = \psi.$

      \item For an arbitrary $\varphi \in (V/U)',$ we have
        $\pi'(\varphi) = \varphi \circ \pi.$ Hence, let us consider
        some $(\pi'(\varphi))(u)$ for $u \in U.$ Then, we get
        $$
        (\pi'(\varphi))(u) = \varphi(\pi(u)) = \varphi(u + U) =
        \varphi(U) = \varphi(0 + U)= 0,
        $$
        since any functional $x \in \mathcal L(V/U, \mathbb F)$ must
        satisfy $\varphi(U) = 0.$ Hence, $\pi'(\varphi) \in U^{0}$
        for any $\varphi \in (V/U)',$ and thus $\text{range}\ \pi' = U^{0}.$

      \item Since $\pi'$ is injective, and $\text{range} \ \pi' =
        U^{0},$ then if we restrict the codomain of $\pi'$ to
        $U^{0},$ it must be an isomorphism between the two spaces.
    \end{enumerate}
\end{enumerate}

\subsection{Reading}
\begin{enumerate}
  \item[13]
    \begin{enumerate}
      \item Consider $\mathcal P(\mathbb F) / U,$ i.e., the set such that
        $$
        \mathcal P(\mathbb F) /U = \{r + U : r \in \mathcal P(\mathbb F)\}.
        $$

        For simplicity, first assume $\deg p > 0.$

        We show that $\dim \mathcal P(\mathbb F) / U = \deg p$ by
        constructing an explicit basis. Construct the set
        $$
        \{1 + U, z + U, z^2 + U, \ldots, z^{m - 1} + U\} \subseteq
        \mathcal P_{m - 1}(\mathbb R).
        $$

        We prove that this set is both linearly independent and spanning.

        \begin{itemize}
          \item
            For linear independence, assume we have
            $$
            \lambda_0(1 + U) + \ldots + \lambda_{m - 1}(z^{m - 1} + U) = 0 + U
            $$
            for coefficients $\lambda_0, \ldots, \lambda_{m - 1} \in
            \mathbb F.$ Then, we get
            $$
            (\lambda_0 + \lambda_1 z + \ldots + \lambda_{m
            -1}z^{m-1}) + U = 0 + U.
            $$
            Hence, we need $\lambda_0 + \lambda_1 z + \ldots +
            \lambda_{m-1}z^{m-1} \in U,$ i.e., we must be able to
            write it as $pq$ for some $q \in \mathcal P(\mathbb F).$
            Assume for the sake of contradiction that $p, q \neq 0.$
            Then, this polynomial has degree $m - 1,$ and $p$ has
            degree $m,$ meaning $pq$ must have degree $\geq m,$ which
            is a contradiction. Thus, $p$ or $q$ is the zero
            polynomial, meaning $$\lambda_0 + \ldots +
            \lambda_{m-1}z^{m -1} = 0,$$ and thus the coefficients
            satisfy $\lambda_0 = \ldots = \lambda_{m-1} = 0.$ This
            means the set must be linearly independent.

          \item For spanning, consider an arbitrary $x \in \mathcal
            P(\mathbb F) / U,$ i.e, $x = r + U$ for some $r \in
            \mathcal  P(\mathbb F).$ Consider cases on $\deg r$:
            \begin{itemize}
              \item If $\deg r \geq m,$ using the division algorithm
                we can write $r = pq + r'$ for some $\deg r'< m,$ and
                $q \in \mathcal P(\mathbb F)$ Then, since $1, z,
                \ldots, z^{m-1}$ is a basis for $\mathcal
                P_{m-1}(\mathbb F),$ we can write
                $$
                r= pq + \sum_{k=0}^{m-1} \lambda_k z^{k},
                $$
                i.e.,
                $$
                x = \left(\sum_{k = 0}^{m -1} \lambda_{k} z^{k} + pq\right) + U.
                $$
                Since $pq \in U,$ we get $$
                x = \sum_{k = 0}^{m-1} \lambda_k (z^{k} + U),
                $$
                meaning $x$ can be written as a linear combination of
                our proposed basis.

              \item Otherwise, if $\deg r < m$ we choose $q = 0$ in
                the above, and still yield $$x = \sum_{k = 0}^{m-1}
                \lambda_k(z^{k} + U)$$ for some coefficients
                $\lambda_k \in \mathbb F.$
            \end{itemize}

            In either case, $x \in \text{span}(\cdot),$ for all $x
            \in \mathcal P(\mathbb F) / U,$ meaning the set is spanning.
        \end{itemize}

        Hence, the set $\{z^{k} + U : 0 \leq k < m, k \in \mathbb
        Z\}$ is a basis on $\mathcal P(\mathbb F) / U.$ Since it has
        size $m = \deg p,$ we conclude that
        $$
        \dim \mathcal P(\mathbb F) / U = \deg p.
        $$

      \item By the above,
        $$
        \{z^{k} + U : 0 \leq k < m, k \in \mathbb Z\}
        $$
        is a basis.
    \end{enumerate}
\end{enumerate}
\end{document}
